\chapter{Mapping the Code onto JAX and the GPU}
\label{ch:gpu}

This chapter is the implementation blueprint: the JAX vocabulary you need,
data layouts, batching strategy, precision policy, and JIT discipline. No
program listings --- but every decision you will face is settled here.
Nothing in it requires a GPU: the same program runs unchanged on a laptop
CPU, where the small cases of Chapter~\ref{ch:verify} take a minute or two.
Develop there; move to an accelerator when the physics is right.

\section{The JAX you need}
\label{sec:jaxprimer}

JAX is NumPy with three additions: it runs on accelerators, it can
differentiate functions, and it can compile them. If you can write NumPy
you can write JAX (\texttt{import jax.numpy as jnp} and use \texttt{jnp}
where you would use \texttt{np}), provided you respect a few rules that
differ from NumPy. The complete list of what this project uses:
{\small
\begin{longtable}{@{}>{\raggedright\arraybackslash}p{0.27\textwidth}p{0.68\textwidth}@{}}
\toprule
JAX feature & What it does, and where you use it\\
\midrule
\endfirsthead
\toprule
JAX feature & What it does, and where you use it (continued)\\
\midrule
\endhead
\bottomrule
\endlastfoot
\texttt{jax.config.update(} \texttt{"jax\_enable\_x64", True)} & Must be the first
statement of the program. Without it JAX silently computes everything in
single precision and none of the reference numbers can be reproduced.\\
\texttt{jnp.fft.fftn}, \texttt{ifftn} with \texttt{axes=} & 3D
transforms over the last three axes of a batch; same normalization as
NumPy (Fourier conventions of Chapter~\ref{ch:solids}).\\
\texttt{jnp.einsum}, \texttt{@} & All projections, Gram matrices and
subspace rotations.\\
\texttt{jnp.linalg.eigh} & Hermitian eigenproblems (complex input is
fine); eigenvalues come back ascending.\\
\texttt{x.at[idx].set(v)}, \texttt{.add(v)} & JAX arrays are immutable:
\texttt{x[idx] = v} is an error. This form returns a modified copy; it
is how coefficients are scattered into the FFT box. Reading
(\texttt{x[idx]}) works as in NumPy. \emph{Pitfall:} the padded entries
of a $\kk$-point's index array must point somewhere; if they point to box
index 0 and you use \texttt{.set}, a padded zero can overwrite the true
$\GG = 0$ coefficient. Use \texttt{.add} (padded coefficients are zero, so
they add nothing).\\
\texttt{jnp.where(cond, a, b)} & Replaces \texttt{if} on array values
(masks, thresholds, soft locking). Both branches are evaluated, so make
both finite: for ``$1/s$ where $s > \tau$, else $0$'' write
\texttt{jnp.where(s > tau, 1/jnp.where(s > tau, s, 1.0), 0.0)} --- the inner
\texttt{where} only makes the discarded branch finite, the outer one still
supplies the zeros (dropping it gives $1$, not $0$, below the threshold:
in Algorithm~\ref{alg:ortho} that would keep dependent directions). Never
write \texttt{jnp.where(s > tau, 1/s, 0.0)} with a zero in \texttt{s}.\\
\texttt{jax.jit(f)} & Compiles \texttt{f} for fixed array shapes and
dtypes. Inside, Python \texttt{if}/\texttt{while} may not depend on array
\emph{values}; loops with a fixed count are simply unrolled.\\
\texttt{jax.vmap(f)} & Turns a function written for one $\kk$-point (or
one grid point) into one for a whole batch. Write Algorithms
\ref{alg:applyh}--\ref{alg:lobpcg} for a single $\kk$ and
\texttt{vmap} them.\\
\texttt{jax.lax.map(f, xs)} & Like \texttt{vmap} but sequential: use it
where a full batch would not fit in memory (density build, large
meshes).\\
\texttt{jax.grad(f)} & Derivative of a scalar function. Gives $\vxc$
from the XC energy density (Chapter~\ref{ch:xc}): \texttt{grad} of
$f(n,\varsigma)$ with respect to each argument, then \texttt{vmap} over
grid points.\\
\end{longtable}}
\noindent Three habits prevent most JAX-specific trouble. First, do all
geometry-time setup (the integer bookkeeping of Algorithm~\ref{alg:basis},
projector tables, Ewald) in ordinary NumPy and convert the finished arrays
with \texttt{jnp.array}; only the numerical kernels need to be JAX.
Second, keep the outer SCF loop, the Fermi-level bisection and the Pulay
solve in plain Python at first, calling jitted kernels from it --- you can
print and inspect everything. Third, the first call of a jitted function
is slow (it compiles) and it recompiles whenever a shape changes; give
every $\kk$-point the same padded shape so that this happens once.

\section{The parallelism inventory}

Plane-wave DFT exposes four nested, independent axes of parallelism:
\begin{center}
\begin{tabular}{@{}llll@{}}
\toprule
Axis & Size (Au, $q{=}11$) & Independence & GPU mapping\\
\midrule
$\kk$-points & $N_k \sim 10^3$ (full $12^3$ mesh)
& fully independent linear algebra & leading batch dim\\
Bands & $N_b \sim 24$ & coupled only via Rayleigh--Ritz & block/GEMM dim\\
Spinor & 2 & coupled by SOC and $\mathbf B_{\mathrm{xc}}$
& folded into band dim\\
$\GG$ / grid & $N_G \sim 2{\times}10^3$, grid $\sim 40^3$
& elementwise / FFT & thread-level (XLA)\\
\bottomrule
\end{tabular}
\end{center}
The design rule: \textbf{never write a loop over any of these axes ---
express every operation as one array program over the full
$[N_k, N_b, 2, N_G^{\max}]$ tensor}, and let vectorized primitives
(\texttt{vmap}, batched FFT, batched GEMM, batched \texttt{eigh}) fill the
device. A $12^3$ mesh $\times$ 24 bands $\times$ 2 spin components $\approx
8\times10^4$ FFTs of size $\sim 36^3$ per Hamiltonian application: enough
to saturate any current accelerator (and too many to hold in memory at
once --- see the memory budget below).

\section{Data layout and static shapes}

\begin{itemize}[leftmargin=2em]
\item \textbf{Wavefunctions} $C$: complex array $[N_k, N_b, 2, N_G^{\max}]$.
Each $\kk$-point's cutoff sphere has a different size $N_G(\kk)$; pad to
$N_G^{\max} = \max_\kk N_G(\kk)$ and carry a boolean mask
$[N_k, N_G^{\max}]$. Masked entries are kept exactly zero (multiply by the
mask after every operation that could pollute them). Padding overhead is a
few percent; the reward is one static shape for everything, which JAX/XLA
requires for compilation and which makes batched FFT/GEMM trivial.
\item \textbf{Sphere $\leftrightarrow$ box maps}: for each $\kk$, a
precomputed integer array of linear FFT-box indices for its $\GG$ sphere.
Pad-and-scatter ($\texttt{zeros(box).at[idx].add(coeffs)}$-style) and
gather are the only ``indexing'' operations in the code.
\item \textbf{Fields} ($n$, $\mathbf m$, potentials): real/complex arrays on
the dense grid $[N_1, N_2, N_3]$ (and their Fourier images). Mixing history:
$[M_{\mathrm{hist}}, 4, N_1, N_2, N_3]$ in Fourier space.
\item \textbf{Projectors}: complex $[N_k, N_{\mathrm{proj}}, N_G^{\max}]$;
coupling blocks $D$: complex
$[N_{\mathrm{proj}} \times 2, N_{\mathrm{proj}} \times 2]$ block-diagonal
(store dense; it is tiny).
\end{itemize}

\section{The seven kernels}

The entire program reduces to seven device kernels; everything else is glue.
\begin{enumerate}[leftmargin=2em]
\item \textbf{Batched zero-pad / gather} between sphere and FFT box.
\item \textbf{Batched 3D FFT / IFFT} over the $[N_k \times N_b \times 2]$
batch. In JAX, \texttt{jnp.fft.fftn} with \texttt{axes=(-3,-2,-1)} transforms
only the last three axes of a batched array such as
$[N_k, N_b, 2, N_1, N_2, N_3]$ and is therefore already batched (without
\texttt{axes=} it would transform every axis); sizes with small prime factors keep cuFFT on fast paths.
\item \textbf{Elementwise field kernels}: kinetic multiply, local-potential
multiply, XC evaluation, Kerker filter, preconditioner, smearing functions.
XLA fuses chains of these into single kernels automatically.
\item \textbf{Batched tall GEMMs}: projections $B^{\dagger}C$ and expansions
$B(DP)$; subspace products $S^{\dagger}(HS)$; block updates $SY$. These are
\texttt{einsum}s over the $[N_k]$ batch.
\item \textbf{Batched small \texttt{eigh}}: the $3N_b \times 3N_b$
Rayleigh--Ritz problem and the $N_b\times N_b$ Gram matrices of
\textsc{Orthonormalize} (Algorithm~\ref{alg:ortho}; four calls per
iteration), batched over $\kk$
(\texttt{vmap(eigh)}). No Cholesky and no QR anywhere: \texttt{eigh} never
fails, which matters because a jitted program cannot branch on a failed
factorization.
\item \textbf{Global reductions}: density accumulation
($\sum_{n\kk} w f |\psi|^2$ --- an einsum-reduction after the IFFT batch),
electron count for Fermi bisection, residual norms, energies.
\item \textbf{Scatter of constants}: structure factors, projector tables,
Ewald --- geometry-time, negligible.
\end{enumerate}

\section{Precision policy}

\begin{itemize}[leftmargin=2em]
\item Run wavefunctions and FFTs in \textbf{complex64} (single precision)
during the bulk of the SCF: half the memory and typically $2$--$4\times$
the speed of double precision (more on consumer GPUs, whose double-precision
throughput is throttled).
\item Single precision has a floor: residuals cannot be pushed much below
$\sim10^{-5}$ and energies carry $\sim10^{-6}$ relative noise. Finish with
a few SCF iterations --- and do all band-structure, $E(V)$ and topology
post-processing --- in \textbf{complex128}. The tight tolerances quoted
elsewhere in this document ($10^{-8}$ and below) refer to double
precision.
\item Accumulate \textbf{energies, the density, and the small Gram and
Rayleigh--Ritz matrices in float64/complex128}. The subspace matrices are
tiny, so promoting them is free and removes the main weakness of
single-precision LOBPCG (loss of orthogonality). In single precision you
will also need to raise the discard thresholds of
Algorithm~\ref{alg:lobpcg} --- as a starting point, to about
$10^{-5}$ (for $W$) and $10^{-3}$ (for $P$): a direction that is linearly
dependent to seven digits carries no information in a seven-digit
arithmetic.
\item \textbf{Write and validate the whole program in double precision
first} (every reference number of Chapter~\ref{ch:verify} is a
double-precision number). Single precision is an optimization for later:
be \emph{explicit} about dtypes at every array creation, keep full
float64 as a debug switch --- energy differences between the two runs
localize precision bugs instantly.
\item Target after the double-precision finish: total energy reproducible
to $\lesssim 10^{-6}$\,Ha/atom; eigenvalues to $10^{-7}$\,Ha.
\end{itemize}

\section{JIT discipline}

\begin{itemize}[leftmargin=2em]
\item \textbf{One jitted SCF step}: a single function
(fields, $X$, history) $\mapsto$ (fields', $X'$, history', scalars),
containing the eigensolver steps, Fermi search, density build, and mixing.
The Python-level SCF loop calls it repeatedly; convergence is checked on the
returned scalars. (A fully on-device \texttt{lax.while\_loop} SCF is
possible but makes debugging miserable --- do it last, if ever.)
\item \textbf{Static shapes everywhere}: fixed $N_b$, fixed LOBPCG
iterations per SCF step, fixed mixing history length (a ring buffer),
fixed bisection iteration count. Convergence-dependent behavior is expressed
through \texttt{where}-masks, not shape changes.
\item \textbf{No host--device sync inside the step}: never branch Python
control flow on device values inside the jitted region; return diagnostics
as arrays.
\item \textbf{Multi-GPU} (when you outgrow one device): shard the leading
$\kk$ axis (\texttt{shard\_map}); the only cross-device
communication per SCF step is (i) an all-reduce of the density and (ii) an
all-reduce of the eigenvalue histogram for the Fermi search --- a few MB.
This near-embarrassing parallelism over $\kk$ is why the algorithmic choices
of Chapter~\ref{ch:algorithms} scale.
\item \textbf{Autodiff opportunities}: XC potentials via \texttt{grad} of the
energy density (Ch.~\ref{ch:xc}); later, forces and stress via
differentiating the total energy with respect to atomic positions and
strain --- the Hellmann--Feynman theorem emerges automatically, provided
you differentiate the energy expression at \emph{fixed}, converged
wavefunctions and occupations (differentiating through a truncated SCF
loop gives something else). Design the
energy as a pure function of (parameters, density, wavefunctions) from day
one.
\end{itemize}

\section{Memory budget (sanity check)}

Full $\Gamma$-centered $12^3$ mesh, $N_b = 24$, $N_G^{\max} \approx 1500$ ($E_{\mathrm{cut}}
= 40$\,Ha; $\approx 2500$ at $60$\,Ha), complex64:
wavefunctions $1728 \times 24 \times 2 \times 1500 \times 8\,\mathrm{B}
\approx 1.0$\,GB. LOBPCG holds six such arrays ($X, W, P$ and their images
under $H$), so budget $\sim 6$\,GB; time-reversal symmetry halves $N_k$
for free (Section~\ref{sec:symmetry}). The FFT boxes are the real constraint: all states at once on a
$36^3$ grid would be $1728\times24\times2\times36^3\times8\,\mathrm{B}
\approx 31$\,GB, so the $\hat H$-apply and the density build \emph{must} be
processed in $\kk$-chunks (a static chunk size, scanned with
\texttt{lax.map}; 64 $\kk$-points per chunk is $\sim 1$\,GB). Fields and
mixing history are megabytes.

\begin{keybox}{Recommended parameters for the first converged gold run}
HGH PBE $q{=}11$; $E_{\mathrm{cut}} = 40$\,Ha (FFT grid $36^3$);
$\Gamma$-centered $12{\times}12{\times}12$ mesh; Methfessel--Paxton
order 1, $\sigma = 0.01$\,Ha; $N_b = 24$ spinor bands; LOBPCG: 8
iterations in the first SCF step, 3 per step afterwards, Teter
preconditioner; Pulay history 8, Kerker $A = 0.8$, $k_{TF} = 1.0$;
convergence $\|R\| < 10^{-7}$ (Eq.~\eqref{eq:resnorm}). Expect about ten
SCF iterations. Then converge upward: $E_{\mathrm{cut}} \to 60$, mesh
$\to 16^3$, and verify observables move by less than your target
accuracy. For debugging, the small case of Chapter~\ref{ch:verify}
(Stage~3, Table~\ref{tab:smallcases}) runs in about a minute on a CPU.
\end{keybox}
